\section{Unallocated capacity in the original microscopic source}
\label{sec:v107-allocation}
A comparison family need not satisfy full envelope balance on every
block. The correct decomposition then retains the uncovered part
as a positive source, rather than discarding the block or replacing
its trace by the comparison envelope.

\begin{theorem}[Fractional comparison capacity]
\label{thm:v107-allocation}
Assume Definition~\ref{def:v107-family}, without
\eqref{eq:v107-balance}. Put
\begin{equation}\label{eq:v107-beta}
 \beta_j=\begin{cases}
       \min\{1,\kappa P_j/C_j\},&C_j>0,\\
       0,&C_j=0.
       \end{cases}
\end{equation}
Then $\sum_j\beta_jT_j$ satisfies
\eqref{eq:v107-finite}--\eqref{eq:v107-height}. At every finite count,
\begin{equation}\label{eq:v107-exact-block-split}
 \sum_jT_j=\sum_j\beta_jT_j+\sum_j(1-\beta_j)T_j,
 \qquad (1-\beta_j)C_j=(C_j-\kappa P_j)_+.
\end{equation}
For a fixed comparison family, these are the largest scalar
fractions in $[0,1]$ for which the envelope-balance certificate holds
block by block.
\end{theorem}
\begin{proof}
Replace $s_j$ by $\beta_js_j$ in
Definition~\ref{def:v107-family}. Its new source capacity is
$\beta_jC_j\le\kappa P_j$; no comparison weight or physical segment
has changed. Theorem~\ref{thm:v107-envelope} applies. The displayed
identities follow from \eqref{eq:v107-beta}. If a scalar fraction
$\gamma_j\in[0,1]$ satisfies $\gamma_jC_j\le\kappa P_j$ and $C_j>0$,
then $\gamma_j\le\beta_j$. This is optimality for the stated
capacity certificate, not an optimality claim over all possible
physical comparison constructions.
\end{proof}

\begin{corollary}[Two endpoints improve one-ended capacity]
\label{cor:v107-monotone}
For any fixed family, replacing $P_j$ by either
$P_j^-:=\int w_jf_-\,d\theta$ or
$P_j^+:=\int w_jf_+\,d\theta$ cannot increase $\beta_j$.
Consequently use of the two-ended envelope cannot enlarge the
unallocated source for that family and budget.
\end{corollary}
\begin{proof}
Pointwise $p_\vee\ge f_-$ and $p_\vee\ge f_+$. Integrate and use monotonicity
of the function $x\mapsto\min\{1,\kappa x/C_j\}$ when $C_j>0$.
Zero-capacity blocks contribute zero under either choice.
\end{proof}

We apply the result to the original physical remainder, not to an
auxiliary source. Fix the outer parameter $B$ and a previously fixed
$0<h_B<\delta(B)$. Let $\widetilde H_B$ denote the original unmatched Borel selector
of Theorem~\ref{thm:v105-original}: it retains the failure source
outside both older recovered families, with their first-membership
convention, and is restricted to $0<\ell<h_B$. In particular
\[
 \mathsf U_B^{\rm micro}(t)=\int
 c^{\varepsilon(B)}(\theta)\one_{\{a(\theta)<t<b(\theta)\}}
                  \widetilde H_B(\Psi(\theta,t))\,d\theta.
\]
All evaluations of this selector remain at the original target.
Choose disjoint angular source sets $I_{B,j}$ and comparison weights
satisfying Definition~\ref{def:v107-family} with $s_j=\one_{I_{B,j}}$.
Put $T_{B,j}=T_j$ with target selector $\widetilde H_B$, and define
\begin{align}
 \mathsf V_B(t)&=\sum_j\beta_{B,j}T_{B,j}(t),\label{eq:v107-V}\\
 \mathsf R_B(t)&=
 \int c^{\varepsilon(B)}(\theta)\one_{\{a<t<b\}}
       \widetilde H_B(\Psi(\theta,t))
       \left(1-\sum_j\beta_{B,j}\one_{I_{B,j}}(\theta)\right)
                       \,d\theta.\label{eq:v107-R}
\end{align}
The expression in parentheses belongs to $[0,1]$. It equals $1$
outside the selected angular source sets and equals $1-\beta_{B,j}$
on the $j$th such set.

\begin{theorem}[Exact residual and original pointwise reduction]
\label{thm:v107-original}
The original finite-count source has the positive decomposition
\begin{equation}\label{eq:v107-micro-split}
 \mathsf U_B^{\rm micro}=\mathsf V_B+\mathsf R_B.
\end{equation}
Suppose the older costs in Theorem~\ref{thm:v105-original} obey its
stated budget. Suppose also that the new comparison family has
$0<r_B<\Delta_B\le8\delta(B)$ and
$N_B\kappa_B\le C\varepsilon(B)^{-1/32}$, uniformly in the
collision count, radius and exact labels. Each of these parameters,
including $h_B$, is fixed before the collision limit. Then
\begin{equation}\label{eq:v107-original-error}
 \mathcal E_M\le\mathcal H_M(\mathsf R_B)+C_MB^{-1/384}.
\end{equation}
Along this schedule the unchanged raw pointwise assertion
\eqref{eq:v76-target} is equivalent to
\begin{equation}\label{eq:v107-final-height}
 \mathcal H_M(\mathsf R_B)\longrightarrow0
                         \quad\text{for every fixed }M.
\end{equation}
\end{theorem}
\begin{proof}
The integral definition and disjointness of the source sets prove
\eqref{eq:v107-micro-split}; for countably many sets use Tonelli.
Theorem~\ref{thm:v107-allocation} and
Remark~\ref{rem:v107-scope} give
$\mathcal H_M(\mathsf V_B)\le C_MB^{-1/384}$.
The preceding exact scalar reduction and subadditivity of the
nonnegative essential-height functional therefore imply
\eqref{eq:v107-original-error}. If \eqref{eq:v107-final-height}
holds, this bound proves the original assertion. Conversely that
assertion implies $\mathcal H_M(\mathsf U_B^{\rm micro})\to0$
by Theorem~\ref{thm:v105-original}; the pointwise order
$0\le\mathsf R_B\le\mathsf U_B^{\rm micro}$ gives necessity.
\end{proof}

The two-ended estimate pays the entire difference between the
comparison envelope and its original positive paired trace, even
when the latter is zero. It does not claim that every physical
source has available comparison capacity. Formula~\eqref{eq:v107-R}
retains all unmatched sheets, including sheets without a usable
actual corridor, and the full unallocated fraction on every
capacity-deficient block. In particular no inverse-residence or
logarithmic moment is silently inserted into the proof.
The positive limit \eqref{eq:v107-final-height} is the remaining
estimate for the organizing raw pointwise theorem; it has not been
established by the comparisons proved in this article.

\begin{remark}[Compatibility with the earlier insertion criteria]
The untruncated inverse-residence criterion and the logarithmic
criterion retained above remain sufficient tests on this smaller
positive remainder whenever their respective hypotheses hold.
Indeed multiplication by the nonnegative residual fraction in
\eqref{eq:v107-R} preserves pointwise domination by the original
weighted source. Neither finiteness of the weighted source nor its
uniform height is inferred merely from this domination.
\end{remark}
